\section{The return on a fixed space of young displacements}

Fix $K\Subset(0,1/4)$, an integer $k=6$, and $a_s>0$.
Choose $a_s<\gamma/16$ and smaller than the fixed exponential reserves
in the curved near comparison. All constants below may depend on $a_s,k,K$.
Let $\mathscr A$ be the real Banach space of $C^k$ matrix functions on
$(0,\infty)$, trace free and with no entry having exactly one azimuthal
index, with norm
\[
 \|A\|_{\mathscr A}
 =\sup_{s>0}e^{-a_ss}s^{-5/2}
       \sum_{j=0}^k|(s\partial_s)^jA(s)|.
\]
This space is independent of $\nu,\epsilon,\mu$. Its elements need not
be symmetric. Choose two fixed smooth cutoffs. The stream cutoff
$\chi_y(s)=\widehat\chi_y(\nu s)$ equals one below $4/(5\nu)$ and
zero above $1/\nu$. The source cutoff $\chi_f$ equals one through
$1/\nu$ and zero above $1.1/\nu$. Put
\[
 A^y=\chi_f A,\qquad
 B_{\mu,\nu}[A]=(A^y)'+\frac\mu s A^y.
\]
Then $\|B_{\mu,\nu}[A]\|_{\mathscr B_{a_s}}
\le C_K\|A\|_{\mathscr A}$, with the source norm of the preceding
section. The matrix cutoff derivative is included in this definition.
Let $\mathfrak T_{\mu,\nu}A$ denote the complete flat stress of the
zero-entrance forced column with Gaussian age force $f_{B_{\mu,\nu}[A]}$.
The actual cylinder train remains in this evolution.

We next define a return map which retains the entire mean-thin kernel.
For an old thin tensor $M$, use its normalized phase mean and set
\[
 E_{ij}(u)=\frac{\chi_o(u)}c
          \int_{\rm cell}M_{ij}(u,q,\eta)\,dq\,d\eta .
\]
For $x'=(\rho\cos\varphi,\rho\sin\varphi)$, $u=\rho^{-2}$,
let $\Sigma_M(x')$ be the Cartesian rotation of the cylindrical tensor
$\rho^{-1}E(u)$, embedded on $x_3=0$. The zero azimuthal row and column
are retained as zero entries; the $rz$ and $zr$ entries are retained.
Define the exact zero-width Newton field by
\[
 V_{\mu,M}(x)=-2\int_0^1t^{2\mu+2}
       \int_{\mathbb R^2}D^3N(x-t(x',0)):\Sigma_M(x')\,dx'\,dt,
 \qquad N(x)=\frac1{4\pi|x|}.
\]
The contraction is the velocity contraction of $\nabla K_0$, with
$K_0=D^2N$ away from the source. Put
\[
 W^{\rm mt}_{\mu,M}(y)=\epsilon^2V_{\mu,M}(\epsilon y),\qquad
 X^{\rm mt}_{\mu,M}
 =\mathcal C_{\chi_y}\bigl(W^{\rm mt}_{\mu,M}
                              -W^{\rm mt}_{\mu+1,M}\bigr),
\]
where the same tensor $M$ is used in the shifted field. The stream
cutoff uses the outward primitive, so the field is fixed without an
additional entrance constant. Define
\begin{equation}\label{source-class:L}
 (\mathcal L_{\mu,\nu}M)(s)
   =D_yX^{\rm mt}_{\mu,M}(R/\sqrt s,0,0)
       \quad\hbox{in the }(r,\phi,z)\hbox{ frame}.
\end{equation}
This is the full physical gradient, including its hoop entry and all
stream-cutoff terms. It is trace free, has the required azimuthal
symmetry, and vanishes for $s\ge1/\nu$. The uncut field is a harmonic
gradient, but the cutoff gradient need not be symmetric.
In particular $\chi_f\mathcal L M=\mathcal L M$ exactly. Thus the
source map applied to $\mathcal L M$ introduces no second cutoff.

The fixed-space flat return is
$\mathcal K^{\rm flat}_{\mu,\nu}
  =\mathcal L_{\mu,\nu}\mathfrak T_{\mu,\nu}$.
Define the unit displacement and real row by
\[
 a_\mu(s)=\frac{s^{5/2}}{\mu+5/2}D_0,
 \quad D_0=\operatorname{diag}(12,-3,-9),
\]
\[
 \Lambda_{\mu,\nu}(A)
 =-\frac{3\pi}{16}\epsilon^3\nu^{5+2\mu}
      \operatorname{Re}\bigl[iN_\nu J(\mu)
                       a_0(B_{\mu,\nu}[A])\bigr].
\]
Both real harmonic coordinates enter the complex amplitude $a_0$;
the matrix force still contains its radial and hoop components.

\begin{proposition}\label{source-class:flat-limit}
Use the direct young-source theorem and the Gaussian force calculation
of the preceding section.
Uniformly for $\mu\in K$,
\[
 \|\mathcal K^{\rm flat}_{\mu,\nu}
                 -a_\mu\Lambda_{\mu,\nu}\|_{\mathscr A\to\mathscr A}
 \le C A_*\left[\nu+\eta_\nu+\omega_{\rm can}(\nu)
      +P_\nu e^{-\gamma'/\sqrt\nu}
      +P_\nu e^{-4a_s/(5\nu)}\right].
\]
The constants multiplying $\nu,\eta_\nu,\omega_{\rm can}$ are fixed,
not viscosity polynomials. Also
$\|a_\mu\|_{\mathscr A}\le C$ and
$\|\Lambda_{\mu,\nu}\|\le C A_*$.
The operators are real and continuous in $\mu$ for each fixed geometry.
On every band where $A_*$ stays bounded, the displayed error tends to zero.
For the unit displacement,
\[
 \Lambda_{\mu,\nu}(a_\mu)
 =2\epsilon^3\nu^{5+2\mu}
       \operatorname{Re}\left[\frac{63\pi i}{128}
                               L(\mu)J(\mu)N_\nu\right]
       +O(A_*P_\nu e^{-b/\nu}).
\]
\end{proposition}
\begin{proof}
Write $M=\mathfrak T_{\mu,\nu}A$, with $\|A\|_{\mathscr A}\le1$.
The uniform source theorem gives, for $j=0,1$,
\[
 \int_0^\infty u^{-3/2-j/2}\frac{\chi_o(u)}c
        \int_{\rm cell}|M(u,q,\eta)|\,dq\,d\eta\,du
       \le C_j\mathcal A_\nu\nu^j.
\]
Its earlier-old part $1.2/\nu<u<u_o=\nu^{-3/2}$ instead has
$\mathcal A_\nu P_\nu e^{-\gamma'/\sqrt\nu}$ as a bound, including
any fixed extra age powers. All tensor components are estimated together.

For the analytical tail $u\ge u_o$, the source radius satisfies
$\rho/r\le\nu^{1/4}<1/2$ on $s\le1/\nu$.
The integrated Cartesian tensor has the form
$\operatorname{diag}(T_h,T_h,T_z)$, with
$T_z-T_h=(\pi/2)I_T(M;u\ge u_o)$. The cross tensor has zero
integrated Cartesian moment. The isotropic part contracts to zero,
since $\partial_i\Delta N=0$ off source. Direct differentiation gives
\[
 D^2\partial_{33}N(r,0,0)
     =(4\pi r^5)^{-1}\operatorname{diag}(-12,3,9).
\]
The negative sign in $V_{\mu,M}$ and the dilation integral therefore give
\[
 D_yW_{\mu,M}^{(0)}
    =\frac{\epsilon^3 I_T(M;u\ge u_o)}{8(\mu+3/2)}s^{5/2}D_0.
\]
This calculation includes the hoop derivative $W_r/|y'|$.
Subtracting the shifted field, with the same tensor, supplies the
additional factor $(\mu+5/2)^{-1}$.

Every point in the mean-value remainder lies at distance at least
$r/2$ from zero. The fifth Newton derivative is bounded by $Cr^{-6}$.
Its radial first moment is the $j=1$ moment above. Thus, for all
$0\le j\le k$,
\[
 |(s\partial_s)^jD_y(W^{\rm mt}_{\mu,M}-W^{\rm mt}_{\mu+1,M}
                     -W^{(0)}_{\mu,M}+W^{(0)}_{\mu+1,M})|
       \le C_k\epsilon^3\mathcal A_\nu\nu s^3.
\]
Receiver conormal derivatives act only on the smooth separated kernel;
no derivative of $M$ is taken. This bound includes all higher source
multipoles, without expanding them one at a time. Since
$\sup_{s>0}s^{1/2}e^{-a_ss}<\infty$, its $\mathscr A$ norm is
at most $C_{a_s,k}A_*\nu$.

The same separated-kernel bounds, on the outward primitive ray,
control the stream cutoff and its fixed conormal derivatives. The
analytical-tail transition contribution is bounded by
$P_\nu A_*e^{-4a_s/(5\nu)}$ in $\mathscr A$.
The uncut leading displacement beyond the cutoff has the same bound.
The earlier-old source has a fixed fractional gap even at
$u=1.2/\nu,s=1/\nu$; its entire exact kernel, including the stream
cutoff, is bounded by $P_\nu A_*e^{-\gamma'/\sqrt\nu}$.
Consequently
\[
 \|\mathcal L_{\mu,\nu}M-a_\mu\ell_\mu(M)\|_{\mathscr A}
 \le C A_*[\nu+P_\nu e^{-\gamma'/\sqrt\nu}
                      +P_\nu e^{-4a_s/(5\nu)}].
\]
The uniform scalar estimate in the preceding section replaces
$\ell_\mu(M)$ by $\Lambda_{\mu,\nu}(A)$ and proves the result.

The row is bounded because its Laplace integrand is bounded by
$C t^{\mu+3/2}e^{-(\gamma-a_s)t}$.
For $A=a_\mu$ its positive amplitude is
$-(21/4)L(\mu)+O(P_\nu e^{-b/\nu})$, with the cutoff derivative
included. This gives the stated coefficient. Continuity in $\mu$
follows by dominated convergence in the Laplace and dilation integrals
and in the exact causal formula. One may use a slightly larger compact
interval within $(0,1/4)$ for the dominating old-age powers.
The finite conormal norms are dominated by the same separated kernels,
so this gives operator-norm continuity, not merely pointwise continuity.
\end{proof}

For later use, the smooth extraction map has a bound on arbitrary
old thin tensors. If
\[
 N_0(M)=\int_{1.2/\nu}^\infty u^{-3/2}
                \frac1c\int_{\rm cell}|M|\,dq\,d\eta\,du,
\]
then the separated fourth-derivative kernel and its outward primitive give
\begin{equation}\label{source-class:ambient-L}
 \|\mathcal L_{\mu,\nu}M\|_{\mathscr A}
       \le C_{K,k,a_s}\epsilon^3N_0(M).
\end{equation}
Indeed each fixed $r\partial_r$ derivative costs the same power $r^{-5}$;
the fractional radial gap is fixed on the whole young--old rectangle.
Conormal derivatives of $\widehat\chi_y(\nu s)$ are uniformly bounded.
In the gain stress space, $N_0(M)\le P_\nu e^{\gamma/\nu}
\|g^{-1}M\|_Z$, hence the norm of $\mathcal L$ is polynomial on a
bounded loop-gain band. This estimate uses only the thin component.
It assigns no thin moment to the broad repair tensor.

The curved loop is $\mathcal L D_N^{-1}M_UR_\mu F_A$, with the
original cutoff to the near region and the exact physical source below. The proposition
does not identify that loop with $\mathcal K^{\rm flat}$. The equation
comparison supplies their difference as an exponentially small operator
on $\mathscr A$, after the factor $e^{1.1a_s/\nu}$ is paid. The remaining
velocity witnesses belong to the mixed ambient block. Thus no polynomial
inverse is applied to the $O(\nu)$ or canonical-modulus terms above.
